\subsection{Geometric constraints}
The preceding examples already contain rates in their governing equations. A
geometric constraint shows why differentiation is sometimes needed to obtain
a differential representation. A point moving on a circle of radius $R$
obeys 
\begin{equation}
f(x,y)=x^2+y^2-R^2=0.  \label{eq:circular_constraint}
\end{equation}
This equation relates positions but contains no motion rates.
Differentiating it along any differentiable admissible path gives 
\begin{equation}
2x\dot x+2y\dot y=0.  \label{eq:circular_differential}
\end{equation}
Thus, with $z=\func{col}(x,y)$, the coefficient field is indeed the Jacobian
of the constraint: 
\begin{equation}
F(z)=D_zf(z)=%
\begin{bmatrix}
2x & 2y%
\end{bmatrix}%
,  \label{eq:circular_jacobian}
\end{equation}
and the differential representation is $F(z)\dot z=0$. The distinction from
the mechanical example is now explicit: $F$ is obtained by differentiating $%
f $, whereas $H$ was read directly from the force balance. Moreover, the
differential field in (\ref{eq:circular_jacobian}) contains no information
about $R$. Zero incompatibility can identify a relationship consistent with
observed motion without identifying the particular level set; an observed
initial position supplies the missing radius.

